


\chapter{Deployment}
A \textbf{Machine Learning (ML) pipeline} is a structured sequence of automated processes that facilitate the end-to-end workflow of developing, training, evaluating, and deploying machine learning models. It ensures efficient data handling, model optimization, and real-time inference by integrating multiple stages in a streamlined, reproducible manner\cite{Panebianco:2018}.
\section{ML Structure}
 The structure of the ML pipeline for gait pattern recognition consists of several stages, each addressing critical tasks necessary for building, training, and deploying the classification model. \\These stages include data collection from IMUs, data preprocessing to clean and standardize sensor readings, feature extraction to identify relevant motion patterns, model training to classify gait patterns, hyperparameter tuning to optimize model performance, model evaluation to validate accuracy, and deployment on the Arduino Nano 33 BLE Sense for real-time classification\cite{Panebianco:2018}.
\subsection{Objectives of MLPipeline}
The primary idea behind our ML pipeline is to develop a reliable, scalable, and efficient system that can accurately classify gait patterns as either normal or Parkinsonian using motion data from IMUs. The pipeline automates the process from data collection to model deployment, enabling real-time classification and immediate feedback. This approach ensures the system can be used effectively for early detection and monitoring of Parkinson's disease in a portable and cost-effective manner.\cite{Patel:2023}
\subsection{ML Pipeline}

\begin{figure}[H]
    \centering
    \begin{tikzpicture}[
        node distance=1cm,
        box/.style={rectangle, draw, text width=6cm, text centered, minimum height=1cm, font=\sffamily,top color =blue!30,bottom color=cyan!30 , rounded corners},
        arrow/.style={->, shorten >=1pt, shorten <=1pt, ultra thick, color=black},
        edge from parent/.style={draw=black, thick}
        ]
        
        % Nodes
        \node[box, top color =blue!30,bottom color=cyan!30] (data_collection) {Data Collection:\\
        - dataloader.py (Loads CSV data)\\
        - normalpatterns.csv (Normal gait patterns)\\
        - parkinsonpatterns.csv (Parkinsonian gait patterns)};

        \node[box, top color =blue!30,bottom color=cyan!30, below=of data_collection] (preprocessing) {Data Preprocessing:\\
        - kalmanfilter.py (Applies smoothing)\\
        - Features: Step\_length, Stride\_length, Step\_time};

        \node[box, top color =blue!30,bottom color=cyan!30, below=of preprocessing] (model_training) {Model Training:\\
        - cnnmodel.py (Defines CNN architecture)\\
        - config.py (Model parameters)\\
        - Training and validation process};

        \node[box, top color =blue!30,bottom color=cyan!30, below=of model_training] (visualization) {Visualization, check accuracy:\\
        - visualization.py (Plots model metrics)};

        \node[box, top color =blue!30,bottom color=cyan!30, below=of visualization] (model_conversion) {Model Conversion:\\
        - tfliteexporter.py (Converts to TFLite format)\\
        - gait\_model.tflite (TFLite model file)};

        \node[box, top color =blue!30,bottom color=cyan!30, below=of model_conversion] (deployment) {Deployment:\\
        - Deploy the model on Arduino Nano 33 BLE Sense\\
        - TFLite Micro Interpreter executes classification\\
        - LED feedback for real-time gait detection};

        % Arrows
        \draw[arrow] (data_collection) -- (preprocessing);
        \draw[arrow] (preprocessing) -- (model_training);
        \draw[arrow] (model_training) -- (visualization);
        \draw[arrow] (visualization) -- (model_conversion);
        \draw[arrow] (model_conversion) -- (deployment);
        
    \end{tikzpicture}
    \caption{ML Pipeline for Gait Pattern Recognition}
    \label{fig:mlpipelinevertical}
\end{figure}


The ML pipeline for our gait pattern classification project includes the following stages:

\begin{enumerate}
    \item \textbf{Data Collection}: Motion data is collected from the inertial measurement unit (IMU) embedded in the Arduino Nano 33 BLE Sense. The dataset includes key gait parameters such as step length, stride length, and step time. Data is captured under controlled conditions for both normal and Parkinsonian gait patterns \cite{ZhangWenwen:2022}.
    
    \item \textbf{Data Preprocessing}: The Kalman Filter is applied to the raw motion data to smooth accelerometer readings, reduce noise, and ensure consistent input. Preprocessed data enhances the reliability and accuracy of the classification model\cite{Kalman:2012}.
    
    \item \textbf{Feature Engineering}: Filtered data is formatted into features such as step length, stride length, and step time, which are suitable for the convolutional neural network (CNN) model. These features allow the model to effectively differentiate between normal and Parkinsonian gait patterns \cite{Panebianco:2018}.
    
    \item \textbf{Model Training}: A CNN is implemented using TensorFlow. The model is trained on 3-dimensional input data (step length, stride length, and step time) to classify gait patterns as either normal or Parkinsonian. Training is performed using a binary cross-entropy loss function and the Adam optimizer to achieve high accuracy\cite{Goodfellow:2016}.
    
    \item \textbf{Model Evaluation}: The trained model is validated using metrics such as accuracy and loss on a separate validation dataset. These metrics help ensure that the model generalizes well to unseen data, providing reliable predictions \cite{Goodfellow:2016}.
    
    \item \textbf{Model Conversion}: After training, the model is converted into TensorFlow Lite format using \texttt{tfliteexporter.py}. This step optimizes the model for deployment on resource-constrained hardware such as microcontrollers\cite{David:2021}.
    
    \item \textbf{Deployment}: The TensorFlow Lite model is deployed on the Arduino Nano 33 BLE Sense. The embedded system processes real-time IMU data, classifies gait patterns, and provides immediate feedback using onboard LED indicators\cite{ArduinoNano33Sense2025}.
\end{enumerate}


\newpage
\subsection{Use in the Project}
The ML pipeline automates the detection and classification of gait patterns. It processes real-time IMU data, classifies the gait using a trained CNN model, and facilitates deployment to the Arduino Nano 33 BLE Sense for real-time analysis.

This ensures that the system is capable of:
\begin{itemize}
    \item \textbf{Data Preprocessing with Kalman Filter}: The pipeline applies the Kalman Filter to smooth sensor data, reducing noise and enhancing feature accuracy \cite{Kalman:2012}.
    \item \textbf{Feature Extraction for Classification}: Extracted features such as step length, stride length, and step time are used to classify gait patterns as normal or Parkinsonian \cite{Panebianco:2018}.
    \item \textbf{Real-Time Classification}: The CNN model, trained on preprocessed data, is deployed on the Arduino Nano 33 BLE Sense to perform real-time classification and provide feedback through an LED indicator \cite{ZhangWenwen:2022}.
    \item \textbf{Ease of Deployment}: The TensorFlow Lite conversion enables efficient deployment of the trained CNN model on resource-constrained hardware, ensuring scalability and portability \cite{Sulaimon:2021}.
\end{itemize}





\section{Steps in the Deployment}

\textbf{Flow Chart :}
The deployment process includes the following:

\begin{figure}[H]
    \centering
    \begin{tikzpicture}[
        node distance=0.5cm,
        box/.style={rectangle, draw, text width=6cm, text centered, minimum height=1cm, font=\sffamily, fill=blue!10, rounded corners},
        arrow/.style={->, shorten >=1pt, shorten <=1pt, ultra thick, color=black},  % Changed to black
        arrowstep/.style={->, shorten >=1pt, shorten <=1pt, ultra thick, color=black},  % Changed to black
        arrowhighlight/.style={->, shorten >=1pt, shorten <=1pt, ultra thick, color=black}  % Changed to black
        ]
        
        % Nodes
        \node[box, top color =blue!30,bottom color=cyan!30] (model_deployment) {Step 1: Model Deployment\\
        - Convert CNN model to TensorFlow Lite format\\
        - Optimize model for resource-constrained devices\\
        - Save as "gait\_model.h"};

        \node[box, top color =blue!30,bottom color=cyan!30, below=of model_deployment] (setup_hardware) {Step 2: Setup Hardware\\
        - Configure IMU sensor\\
        - Define RGB LED pins\\
        - Initialize serial communication};

        \node[box, top color =blue!30,bottom color=cyan!30, below=of setup_hardware] (load_model) {Step 3: Load  Model to hardware\\
        - Load "gait\_model.h" in the TensorFlow Lite Micro environment\\
        - Allocate tensors};

        \node[box, top color =blue!30,bottom color=cyan!30, below=of load_model] (capture_data) {Step 4: Capture Data\\
        - Collect real-time IMU data (X, Y, Z axes)\\
        - Compute acceleration magnitude};

        \node[box, top color =blue!30,bottom color=cyan!30, below=of capture_data] (step_detection) {Step 5: Detect Steps\\
        - Check threshold (e.g., acceleration > 1.2)\\
        - Enforce minimum delay between steps};

        \node[box, top color =blue!30,bottom color=cyan!30, below=of step_detection] (run_inference) {Step 6: Run Inference\\
        - Pass data to TensorFlow Lite model\\
        - Invoke classification};

        \node[box, top color =blue!30,bottom color=cyan!30, below=of run_inference] (feedback) {Step 7: Visual Feedback\\
        - Red LED for normal gait\\
        - Blue LED for Parkinsonian gait\\
        - Log results in serial monitor};

        % Arrows
        \draw[arrow] (model_deployment) -- (setup_hardware);
        \draw[arrow] (setup_hardware) -- (load_model);
        \draw[arrow] (load_model) -- (capture_data);
        \draw[arrow] (capture_data) -- (step_detection);
        \draw[arrow] (step_detection) -- (run_inference);
        \draw[arrow] (run_inference) -- (feedback);
        
    \end{tikzpicture}
    \caption{Deployment Process for Gait Pattern Classification}
    \label{fig:deploymentflowchart}
\end{figure}



\subsection{Description}
The deployment process for gait pattern classification includes the following steps:

\begin{enumerate}
    \item \textbf{Model Deployment}: The trained CNN model is converted into TensorFlow Lite format, making it lightweight and optimized for deployment on the Arduino Nano 33 BLE Sense. This step ensures efficient real-time processing of motion data. The model is then saved as "gait\_model.h" for integration into the Arduino project \cite{Sulaimon:2021}.
    
    \item \textbf{Setup Hardware}: The Arduino Nano 33 BLE Sense is configured for real-time data collection and classification. This includes initializing the IMU sensor for motion data, defining RGB LED pins for feedback, and setting up serial communication for debugging and logging \cite{ArduinoNano33Sense2025}.
    
    \item \textbf{Load TensorFlow Lite Model}: The TensorFlow Lite model is loaded onto the Arduino, validated for schema compatibility, and memory buffers are allocated for input, output, and intermediate computations.\cite{David:2021}.
    
    \item \textbf{Capture Data}: The IMU sensor continuously captures acceleration data from the X, Y, and Z axes. The total acceleration magnitude is computed as the Euclidean norm of the three components, representing the intensity of motion\cite{Panebianco:2018}.
    
    \item \textbf{Detect Steps}: A threshold-based algorithm filters noise and detects valid steps by checking if the acceleration exceeds a predefined threshold (e.g., 1.2). A delay is enforced between consecutive detections to prevent false positives caused by rapid motion\cite{Panebianco:2018}.
    
    \item \textbf{Run Inference}: The acceleration data is fed into the TensorFlow Lite model for classification. The model predicts the likelihood of the detected gait being Parkinsonian or normal\cite{Viswanatha:2022}.
    
    \item \textbf{Visual Feedback}: The classification result is communicated via the RGB LEDs on the Arduino Nano 33 BLE Sense. A red LED indicates normal gait, while a blue LED signals Parkinsonian gait. The step count and classification results are also logged to the serial monitor for real-time monitoring\cite{ArduinoNano33Sense2025}.
\end{enumerate}


\subsection{Configuration}

Configuration is a critical step in setting up the environment and parameters required for implementing the gait pattern recognition pipeline. It ensures smooth interaction between various components and accurate performance throughout the pipeline. The configuration process includes:

\begin{itemize}
    \item \textbf{Data Storage}: 
    Organize raw motion data collected from the Arduino Nano 33 BLE Sense IMU into labeled datasets for training, validation, and testing. Use a structured directory layout to maintain separate folders for different data categories and ensure proper file naming conventions. Each dataset should include parameters such as step length, stride length, and step time \cite{ZhangWenwen:2022}.

    \item \textbf{Kalman Filter Parameters}: 
    Configure the Kalman Filter to smooth and denoise the raw motion data. This involves setting the initial state, transition matrices, process noise covariance, and observation covariance values. Fine-tune these parameters to balance noise reduction and preservation of meaningful gait features \cite{Kalman:2012}.

    \item \textbf{CNN Model Parameters}: 
    Define the architecture and hyperparameters of the convolutional neural network (CNN) used for gait classification. Key parameters include:
    \begin{itemize}
        \item Input shape: (3, 1) to accommodate features (step length, stride length, step time).
        \item Convolutional layers: Number of layers and filter sizes \cite{Panebianco:2018}.
        \item Activation functions: Use ReLU for hidden layers and sigmoid for the output layer.
        \item Hyperparameters: Learning rate, batch size, and number of epochs\cite{Panebianco:2018}.
    \end{itemize}

    \item \textbf{Deployment Environment}: 
    Prepare the Arduino Nano 33 BLE Sense environment for model deployment. This includes:
    \begin{itemize}
        \item Loading the TensorFlow Lite model into the microcontroller \cite{Sulaimon:2021}.
        \item Allocating memory for processing the model and real-time IMU data \cite{ArduinoNano33Sense2025}.
        \item Configuring peripherals such as LEDs to provide feedback for gait classification results (e.g., blue for Parkinsonian gait and red for normal gait) \cite{ArduinoNano33Sense2025}.
    \end{itemize}
\end{itemize}
\newpage


\section{TensorFlow Lite: Description}
LiteRT (short for Lite Runtime), formerly known as TensorFlow Lite\cite{GoogleLiteRT:2023}, is Google's high-performance runtime for on-device AI, tailored to environments with limited computational resources. TFLite enables the deployment of machine learning models on a wide range of devices, from microcontrollers to smartphones and embedded systems. \\

In this project, TensorFlow Lite is crucial for converting a fully trained Convolutional Neural Network (CNN) model into a \texttt{.tflite} format, which significantly optimizes the model for low-resource environments. This conversion facilitates deploying the model on the Arduino Nano 33 BLE Sense, enabling real-time gait pattern classification.\\

For this specific application, the model runs within the TensorFlow Lite Micro (TFLM) environment,a variant of TensorFlow Lite optimized for microcontrollers. TFLM is specifically designed to operate with minimal memory footprint, which aligns perfectly with the constraints of microcontroller units (MCUs) like the Arduino Nano 33 BLE Sense.\\

\textbf{In this project, we use the term ``TensorFlow Lite'' to refer to the broader capabilities of the framework, including its microcontroller-specific adaptations, as it was the well-established version at the time of deployment.} We deliberately avoid the term ``LiteRT'' to prevent confusion with other libraries and maintain clarity in documentation. This decision is based on the fact that, during the deployment stage of this project, TensorFlow Lite underwent a rebranding process to ``LiteRT.'' However, we relied on the well-established TensorFlow Lite library, which has been specifically developed for similar tasks and rigorously tested for our use case. TensorFlow Lite Micro directly supports TensorFlow Lite’s \texttt{.tflite} models without requiring additional runtime environments, making it a reliable and practical choice for our implementation~\cite{GoogleLiteRT:2023}. \\


\subsection{Key Features of TensorFlow Lite}

\begin{itemize}
    \item \textbf{Optimized for On-Device Machine Learning:} Ensures minimal latency, privacy (data remains on the device), and reduced power consumption, crucial for embedded applications\cite{GoogleLiteRT:2023}.
    \item \textbf{Flexible Model Conversion:} Supports model conversion from TensorFlow, PyTorch, and JAX into the lightweight FlatBuffers format (.tflite), which is optimized for edge devices \cite{GoogleLiteRT:2023}.
    \item \textbf{Efficient Inference:} Provides tools for model quantization and leverages hardware acceleration like GPUs or Core ML for high-performance inference.
    \item \textbf{Multi-Platform Support:} Compatible with Android, iOS, Linux, and microcontroller platforms\cite{GoogleLiteRT:2023}.
\end{itemize}

\subsection{TensorFlow Lite: Code Example}
The trained CNN model is converted into the TensorFlow Lite format using Python. Below is the code snippet demonstrating this process:

\lstset{style=pythonstyle}
\lstinputlisting[caption={TensorFlow Lite: Code Example}, label={lst:tensorflow light convert}]{functions/deployment/tensorflowlightconvert.py
}


\subsection{TensorFlow Lite: Use in the Project}
In the project Python script carryout  the TensorFlow Lite conversion via TFLiteConverter, illustrating a crucial step where the model is prepared for deployment by stripping down unnecessary complexities and adapting the model for efficient execution.In this project, TensorFlow Lite is essential for:
\begin{itemize}
    \item \textbf{Model Conversion:} Transforming the trained CNN model into the lightweight .tflite format for deployment.Reduces the precision of the model’s numerical parameters, significantly reducing the model size and often speeding up inference with minimal impact on accuracy\cite{GoogleLiteRT:2023}.
    \item \textbf{Real-Time Optimization:} Enabling efficient inference of motion data on the Arduino Nano 33 BLE Sence.Implements various optimizations that streamline the model, making it feasible for real-time applications on mobile devices and embedded systems\cite{ArduinoNano33Sense2025}.
   
\end{itemize}

\subsection{TensorFlow Lite Micro: Description}
TensorFlow Lite Micro is an extension of TensorFlow Lite specifically designed for microcontrollers. It allows resource-efficient deployment of ML models, such as the gait classification model, on devices with constrained memory and processing power, including the Arduino Nano 33 BLE Sense. 

\textbf{Key Features of TensorFlow Lite Micro:}
\begin{itemize}
    \item \textbf{Low Resource Utilization:} Tailored for microcontrollers with limited computational and memory resources.TensorFlow Lite Micro is designed to run with minimal memory footprint, using techniques like static memory allocation to fit within the few kilobytes of RAM available on microcontrollers.\cite{GoogleLiteRT:2023}.
    \item \textbf{Seamless Integration:} Facilitates real-time inference through direct integration into Arduino IDE and embedded environments. It provides support for the core set of operators required for inference, which allows executing complex machine learning models on a tiny device.\cite{GoogleLiteRT:2023}.
    \item \textbf{Optimized Execution:} Leverages hardware-specific optimizations to execute ML models efficiently on microcontrollers \cite{GoogleLiteRT:2023}.
\end{itemize}

\subsection{TensorFlow Lite Micro: Code Example}
The TensorFlow Lite Micro runtime is used on the Arduino Nano 33 BLE Sense to load the .tflite model and perform inference. Below is a sample implementation:



\lstset{style=pythonstyle}
\lstinputlisting[caption={TensorFlow Lite Micro: Code Example}, label={lst:tensorflow light micro convert}]{functions/deployment/tensorflowlightmicroconvert.py
}

\subsection{TensorFlow Lite Micro: Use in the Project}
TensorFlow Lite Micro is integral to the real-time deployment of the gait classification model. Key functionalities include:
\begin{itemize}
    \item \textbf{Model Deployment:} Loading the .tflite model on the Arduino Nano 33 BLE Sense using TensorFlow Lite Micro\cite{GoogleLiteRT:2023}.
    \item \textbf{Real-Time Classification:} Processing live IMU data from the Arduino's LSM9DS1 sensor to classify gait patterns.
    \item \textbf{User Feedback:} Using LED indicators to provide immediate visual feedback for gait abnormalities.
\end{itemize}
\


\subsection{Use in the Project}
The ML pipeline automates the detection and classification of gait patterns. Its design supports real-time processing and deployment on embedded systems, offering practical solutions for gait analysis.Integrating TensorFlow Lite Micro with Arduino involves setting up the TensorFlow Lite Micro interpreter, loading the model, and using it to make predictions based on real-time sensor data. The code snippet demonstrates initializing the IMU sensor, processing acceleration data, and using the TensorFlow Lite Micro model to classify gait patterns. The process is efficient enough to run within the tight memory and processing constraints of the Arduino platform, showcasing the practical application of machine learning in embedded systems.

This ensures the system is capable of:
\begin{itemize}
    \item \textbf{Data Preprocessing}: The Kalman Filter smooths noisy data, enhancing the reliability of features extracted for classification \cite{Kalman:2012}.
    \item \textbf{Feature Engineering}: Features like step length, stride length, and step time are extracted to classify gait patterns as normal or Parkinsonian \cite{Panebianco:2018}.
    \item \textbf{Efficient Deployment}: The TensorFlow Lite model enables real-time execution on devices with limited resources, such as the Arduino Nano 33 BLE Sense \cite{Sulaimon:2021}.
    \item \textbf{Real-Time Feedback}: The embedded system provides visual feedback through onboard LEDs, offering an interactive user experience\cite{ArduinoNano33Sense2025}.
\end{itemize}


\section{Program optimization}
\subsection{Parameter Handling}
Key parameters in the ML pipeline allow fine-tuning for performance and deployment efficiency. These include:
\begin{itemize}
    \item \textbf{Kalman Filter}: Transition matrices, observation covariance, and noise covariance for data smoothing, which improves signal accuracy by reducing noise\cite{Kalman:2012}.
    \item \textbf{CNN Hyperparameters}: Batch size, learning rate, convolutional layers, and filter size are crucial for optimizing training and improving the model's ability to classify gait patterns\cite{Goodfellow:2016}.
    \item \textbf{Deployment Thresholds}: Acceleration magnitude and step timing thresholds ensure accurate real-time gait detection by reducing false positives and negatives\cite{Zhang:2022}.
    \item \textbf{Hardware Configuration}: TensorFlow Lite settings and pin assignments for IMU sensors enable efficient integration with the Arduino Nano 33 BLE Sense for real-time deployment\cite{ArduinoNano33Sense2025}.
\end{itemize}

\subsection{Error Handling}
Robust error handling ensures the pipeline can address unexpected scenarios gracefully:
\begin{itemize}
    \item \textbf{Data Validation}: Validates datasets for missing values and column consistency to prevent crashes. This ensures the pipeline only processes well-structured data, reducing runtime errors \cite{Pandas:2024}:
    \\lstset{style=pythonstyle}
\lstinputlisting[caption={Error handling: Code Example}, label={lst:errorhandling}]{functions/deployment/errorhandling.py
}

    \item \textbf{Training Failures}: Logs and handles errors during model training, such as data shape mismatches, ensuring the model is trained with valid inputs and preventing wasted computation \cite{TensorFlow:2024}:
      
\lstinputlisting[caption={Train error handling: Code Example}, label={lst:trainerrorhandling}]{functions/deployment/trainerror.py
}

    \item \textbf{Deployment Errors}: Ensures the embedded system is prepared for real-time classification. This step checks hardware readiness, such as IMU initialization, to prevent runtime failures during deployment \cite{ArduinoNano33Sense2025}:
   \lstinputlisting[caption={Deployment error handling: Code Example}, label={lst:deploymenterrorhandling}]{functions/deployment/deploymentfaults.py
}
\end{itemize}

\subsection{Message Handling}
Effective logging provides traceability during pipeline execution:
\lstinputlisting[caption={Message handling: Code Example}, label={lst:messagehandling}]{functions/deployment/messagehandling.py
}
This ensures transparency, enabling developers to monitor the system's behavior during both development and deployment.






\section{Structure of the Program}

\subsection{Project Directory }
The project is organized as follows to streamline the implementation and deployment of the gait classification pipeline:



\lstset{style=pythonstyle}
\lstinputlisting[caption={Project sctructure}, label={lst:projectstructuret}]{functions/deployment/projectstructure.py
}

The program has a modular structure with each module handling a specific aspect of the pipeline:

\begin{enumerate}
    \item \textbf{Main Pipeline (\texttt{main.py})}: The entry point of the program that coordinates the execution of all components, including data loading, preprocessing, model training, and deployment.
    \item \textbf{Logger Setup (\texttt{logger.py})}: Configures logging for real-time feedback during execution.
    \item \textbf{Data Loading (\texttt{dataloader.py})}: Handles data ingestion and validation to ensure it meets the requirements of the pipeline.
    \item \textbf{Data Preprocessing (\texttt{kalmanfilter.py})}: Implements the Kalman Filter for noise reduction and smoothing of sensor data.
    \item \textbf{Model Definition and Training (\texttt{cnnmodel.py})}: Defines and trains the convolutional neural network (CNN) used for classification.
    \item \textbf{Model Export (\texttt{tfliteexporter.py})}: Converts the trained model into TensorFlow Lite format for deployment on embedded systems.
    \item \textbf{Visualization (\texttt{visualization.py})}: Provides tools for visualizing model performance, including accuracy and loss curves.
    \item \textbf{Configuration (\texttt{config.py})}: Centralizes all configurable parameters, such as CNN hyperparameters and Kalman Filter settings.
\end{enumerate}

\subsection{How the Program Works}

The program operates in the following sequence:
\begin{enumerate}
    \item \textbf{Data Collection}: Gait data is collected from IMUs and stored in CSV files (\texttt{normalpatterns.csv}, \texttt{parkinsonpatterns.csv}).
    \item \textbf{Data Preprocessing}: The Kalman Filter is applied to clean and smooth noisy sensor data.
    \item \textbf{Feature Engineering}: Relevant gait metrics (step length, stride length, step time) are extracted and formatted for the CNN.
    \item \textbf{Model Training}: The CNN is trained on the preprocessed data to classify gait patterns as normal or Parkinsonian.
    \item \textbf{Model Export and Deployment}: The trained model is converted to TensorFlow Lite format and deployed on the Arduino Nano 33 BLE Sense.
    \item \textbf{Real-Time Classification}: The deployed model classifies gait patterns in real time, providing visual feedback through LED indicators.
\end{enumerate}

\section{Modules of program}

\subsection{\texttt{main.py}}
[\href{run:../Code/MicroPython/project code/main.py}{Click to export main.py code}] 
\begin{itemize}
    \item Serves as the central controller for the pipeline.
    \item Coordinates tasks such as data loading, preprocessing, model training, and deployment.
    \item Integrates all other modules into a seamless workflow.
\end{itemize}

\subsection{\texttt{logger.py}}
[\href{run:../Code/MicroPython/project code/src/logger.py}{Click to export logger.py code}] 
\begin{itemize}
    \item Configures logging for real-time debugging and monitoring.
    \item Logs all critical steps in the pipeline, including errors and warnings.
\end{itemize}

\subsection{\texttt{dataloader.py}}
[\href{run:../Code/MicroPython/project code/src/dataloader.py}{Click to export dataloader.py code}] 
\begin{itemize}
    \item Loads datasets (\texttt{normalpatterns.csv}, \texttt{parkinsonpatterns.csv}).
    \item Validates data to ensure required columns (\texttt{Step\_length}, \texttt{Stride\_length}, \texttt{Step\_time}) are present.
    \item Raises errors for missing or invalid data.
\end{itemize}

\subsection{\texttt{kalmanfilter.py}}
[\href{run:../Code/MicroPython/project code/src/kalmanfilter.py}{Click to export kalmanfilter.py code}] 
\begin{itemize}
    \item Applies the Kalman Filter to smooth noisy IMU data.
    \item Ensures clean and reliable input data for the CNN.
    \item Enhances model accuracy by reducing noise in gait metrics.
\end{itemize}

\subsection{\texttt{cnnmodel.py}}
[\href{run:../Code/MicroPython/project code/src/cnnmodel.py}{Click to export cnnmodel.py code}] 
\begin{itemize}
    \item Defines the architecture of the CNN:
    \begin{itemize}
        \item Input Layer: Accepts 3-dimensional motion data.
        \item Convolutional Layers: Extract spatial features from the data.
        \item Dense Layers: Perform binary classification.
    \end{itemize}
    \item Trains the model using an Adam optimizer and binary cross-entropy loss function.
    \item Outputs performance metrics such as accuracy and loss.
\end{itemize}

\subsection{\texttt{tfliteexporter.py}}
[\href{run:../Code/MicroPython/project code/src/tfliteexporter.py}{Click to export tfliteexporter.py code}] 
\begin{itemize}
    \item Converts the trained CNN model to TensorFlow Lite format (\texttt{.tflite}).
    \item Generates a \texttt{.h} header file for deployment on Arduino IDE.
\end{itemize}

\subsection{\texttt{visualization.py}}
[\href{run:../Code/MicroPython/project code/src/visualization.py}{Click to export visualization.py code}] 
\begin{itemize}
    \item Visualizes model performance, including training and validation accuracy/loss curves.
    \item Provides insights into model behavior and areas for improvement.
\end{itemize}

\subsection{\texttt{config.py}}
[\href{run:../Code/MicroPython/project code/src/config.py}{Click to export config.py code}] 
\begin{itemize}
    \item Centralizes all configurable parameters.
    \item Includes CNN hyperparameters (e.g., learning rate, batch size) and Kalman Filter settings.
    \item Enables easy experimentation and fine-tuning of the pipeline.
\end{itemize}

\section{How Each Code Contributes to the Objectives}

The program's objective is to achieve real-time classification of gait patterns. Each module contributes as follows:
\begin{itemize}
    \item \textbf{Data Preparation (\texttt{dataloader.py}, \texttt{kalmanfilter.py})}: Ensures clean and reliable input data for the CNN.
    \item \textbf{Model Development (\texttt{cnnmodel.py})}: Builds a CNN capable of learning and classifying gait patterns.
    \item \textbf{Model Export (\texttt{tfliteexporter.py})}: Prepares the trained model for deployment on resource-constrained hardware.
    \item \textbf{Visualization and Logging (\texttt{visualization.py}, \texttt{logger.py})}: Tracks performance and provides feedback for debugging.
    \item \textbf{Real-Time Deployment (\texttt{main.py})}: Integrates all components for seamless real-time classification.
\end{itemize}

\section{Program Workflow}

\subsection{End-to-End Process}
\begin{enumerate}
    \item Load and validate gait datasets.
    \item Preprocess raw IMU data using the Kalman Filter.
    \item Extract relevant features for training the CNN.
    \item Train the CNN to classify gait patterns as normal or Parkinsonian.
    \item Convert the trained model to TensorFlow Lite format for deployment.
    \item Deploy the model on the Arduino Nano 33 BLE Sense for real-time classification.
    \item Receive real-time visual feedback using LED indicators.
\end{enumerate}

\subsection{Program Objectives Achieved}
\begin{itemize}
    \item Accurate classification of gait patterns in real time.
    \item Robust pipeline design with modular components.
    \item Efficient deployment on resource-constrained hardware.
    \item Practical solution for Parkinson's disease monitoring and detection.
\end{itemize}










